\documentclass[10pt,a4paper]{amsart}
\usepackage[margin=19mm]{geometry}
\usepackage{amsmath,amssymb,mathtools}
\usepackage[scr=rsfs,cal=euler]{mathalfa}
\usepackage{xfrac}
\usepackage[unicode,hidelinks,bookmarks=false]{hyperref}
\newcommand{\ie}{i.e.\ }
\newcommand{\cf}{cf.\ }
\newcommand{\resp}{respectively\ }
\newcommand{\Subsection}{\subsection}
\providecommand{\nobreakdash}{\nobreak\hbox{-}}
\setlength{\emergencystretch}{2em}
\multlinegap0pt
\usepackage{bm}
\usepackage{bbm}
\usepackage{dsfont}
\usepackage{xspace}
\usepackage{cancel}
\usepackage{mathtools}
\usepackage{amsthm}
\makeatletter
\@ifundefined{theorem}{\newtheorem{theorem}{Theorem}}{}
\@ifundefined{lemma}{\newtheorem{lemma}[theorem]{Lemma}}{}
\makeatother
\title[Geometric Classification of Biased Quantum Capacity via Harm]{Geometric Classification of Biased Quantum Capacity via Harmonic Translation}
\author{Eliseo Sarmiento Rosales, Egor Maximenko, Dionisio Manuel Tun Molina, Juan Carlos Jimenez Cervantes, Jose Alberto Guzman Vega, Rodrigo Leon Morales}
\date{Source: arXiv:2603.22336v1 (CC BY 4.0)}
\begin{document}
\maketitle

\noindent\emph{Source and licence.} Geometric Classification of Biased Quantum Capacity via Harmonic Translation. Version arXiv:2603.22336v1 (CC BY 4.0), distributed under the Creative Commons Attribution 4.0 International licence, as recorded in the arXiv licence metadata for this exact version and shown on the abstract page https://arxiv.org/abs/2603.22336v1. Full rights notice: licenses/arxiv-2603.22336-CC-BY-4.0.txt.\par\smallskip

\noindent\textit{Editorial scope.} Counted unit: the Lemma whose source label is \texttt{lem:catq\_cayley} in the article (source0.tex lines 1664--1686, source label \texttt{lem:catq\_cayley}), delivered together with the article's own complete original proof of it (source0.tex lines 1688--1731, one \texttt{proof} environment of 44 lines). No mathematics is written, repaired or re-derived: statement and proof are the article's own text line for line. Required context retained uncounted: none. Every definition, display and label of those lines is retained unchanged. Omitted from this package and not redistributed: 1--1663, 1687--1687, 1732--2081. Nothing inside a retained excerpt is omitted or altered: every retained body is delivered exactly as the article's lines are written, so no omission receipt of any kind accompanies this package. Citations: the retained excerpts carry no citation command at all, so no bibliography entry of the article is transcribed and no reference is redistributed. Reference adaptation: none was needed and none was made; every reference inside the delivered unit resolves inside it, so no unresolved reference or citation is produced. Printed slips kept literal: no printed word, symbol, hypothesis or number of the counted statement or of its proof was altered. Layout and class: the article's own class is \texttt{quantumarticle} with packages including tikz, amsmath, amssymb, amsthm, mathtools, bbm, braket, natbib, hyperref; the wrapper is the lane's pinned \texttt{amsart} editorial wrapper at 10pt on A4, which restates the theorem-like environments the retained text needs and adds the article's own macro definitions that the retained text needs (none), with extra wrapper packages (bm, bbm, dsfont, xspace, cancel, mathtools). The delivered numbering is the unit's own. Assets: no retained span contains a figure environment, a graphics-inclusion command, a PDF-page inclusion, a table or a TikZ picture; no image file is copied into this package, the package declares no asset dependency and no image file is redistributed with it. Licence: version arXiv:2603.22336v1 of this article is distributed under the Creative Commons Attribution 4.0 International licence, as recorded in the arXiv licence metadata for this exact version and shown on the abstract page https://arxiv.org/abs/2603.22336v1. A full rights notice is delivered at licenses/arxiv-2603.22336-CC-BY-4.0.txt. Characters: every retained excerpt is pure ASCII, so the delivered engine, XeLaTeX with its default Latin Modern text font, reports no missing character.
\begin{lemma}[Cayley-graph structure under uniform cat-qubit phase noise]
\label{lem:catq_cayley}
Consider an array of $n$ cat qubits in the strongly biased regime
where phase-flip errors dominate and bit-flip errors are negligible.
Let $t\ge 1$ and define the admissible error set
\[
   \Omega \;=\; E_t
     \;=\; \{e\in\mathbb{F}_2^n : \mathrm{wt}(e)\le t\},
\]
representing phase errors affecting at most $t$ qubits. Then:
\begin{enumerate}
\item[\textup{(i)}] The difference set is
  $\Omega-\Omega = E_{2t}$.
\item[\textup{(ii)}] The additive Cayley graph
  $G=\mathrm{Cay}\!\bigl(\mathbb{F}_2^n,\,E_{2t}\setminus\{0\}\bigr)$
  has vertex set $\mathbb{F}_2^n$ with $x\sim y$ if and only if
  $0<d_H(x,y)\le 2t$.
\item[\textup{(iii)}] Independent sets in $G$ are precisely binary
  codes of minimum distance at least $2t+1$.
\item[\textup{(iv)}] The independence number satisfies
  $\alpha(G) = A_2(n,2t+1)$.
\end{enumerate}
\end{lemma}

\begin{proof}
\textbf{(i)}
By definition,
\[
\begin{aligned}
\Omega-\Omega
   &= \{\,e\oplus e' \mid e,e'\in\mathbb{F}_2^n, \\
   &\qquad \mathrm{wt}(e)\le t,\;
            \mathrm{wt}(e')\le t \,\},
\end{aligned}
\]

using $-e'=e'$ in $\mathbb{F}_2^n$.

\emph{Inclusion $\Omega-\Omega\subseteq E_{2t}$.}
For any $e,e'\in E_t$ the triangle inequality gives
$\mathrm{wt}(e\oplus e') \le \mathrm{wt}(e)+\mathrm{wt}(e') \le 2t$,
so $e\oplus e'\in E_{2t}$.

\emph{Inclusion $E_{2t}\subseteq\Omega-\Omega$.}
Let $v\in E_{2t}$, so $\mathrm{wt}(v)\le 2t$. Write
$v=\mathbf{1}_I$ where $I=\mathrm{supp}(v)\subseteq[n]$ with
$|I|\le 2t$. Partition $I$ into disjoint subsets $I_1,I_2$ with
$|I_1|,|I_2|\le t$, set $e=\mathbf{1}_{I_1}$ and
$e'=\mathbf{1}_{I_2}$. Then $e,e'\in E_t$ and
$e\oplus e' = \mathbf{1}_{I_1\cup I_2} = v$.

\textbf{(ii)}
The Cayley graph $G=\mathrm{Cay}(\mathbb{F}_2^n,\,E_{2t}\setminus\{0\})$
has edge set
$\{x,y\}\in E(G)\Leftrightarrow x\oplus y\in E_{2t}\setminus\{0\}
\Leftrightarrow 0<d_H(x,y)\le 2t$,
which is the stated claim.

\textbf{(iii)}
A subset $S\subseteq\mathbb{F}_2^n$ is independent in $G$ if and
only if $d_H(x,y)\ge 2t+1$ for all distinct $x,y\in S$, which is
exactly the defining property of a binary code of minimum distance
at least $2t+1$.

\textbf{(iv)}
The independence number $\alpha(G)$ equals the maximum cardinality
of an independent set in $G$. By (iii), this equals $A_2(n,2t+1)$.
\end{proof}

\end{document}
